\documentclass[12pt, a4paper]{article}
\usepackage{../../../myconfig} % Gọi file cấu hình từ ngoài gốc vào

% ==========================================
% BẮT ĐẦU NỘI DUNG TÀI LIỆU
% ==========================================
\begin{document}

% Truyền 3 tham số: {Nhãn cam} {Chữ to chính giữa} {Chữ ở góc phải Header}
\titlebox{Chủ đề: ĐGNL}{Tính đơn điệu của hàm số}{ĐGNL: Tính đơn điệu}

\textit{\textbf{Dạng 1: Tìm m để hàm số đơn điệu trên các khoảng của nó}}

\drawdashlines{25}
\newpage

% --- CÂU 1 ---
\questionLinesManual{30}{1}{Có bao nhiêu giá trị nguyên của tham số \(m\) sao cho hàm số \(f(x) = \dfrac{1}{3}x^3 + mx^2 + 4x + 3\) đồng biến trên \(\mathbb{R}\).}{%
    \begin{tasks}(4)
        \task 5
        \task 4
        \task 3
        \task 2
    \end{tasks}
}

% --- CÂU 2 ---
\questionLinesManual{30}{2}{Cho hàm số \(y = -x^3 - mx^2 + (4m + 9)x + 5\), với \(m\) là tham số. Hỏi có bao nhiêu giá trị nguyên của \(m\) để hàm số nghịch biến trên khoảng \((-\infty; +\infty)\)}{%
    \begin{tasks}(4)
        \task 5
        \task 4
        \task 6
        \task 7
    \end{tasks}
}

% --- CÂU 3 ---
\questionLinesManual{30}{3}{Cho hàm số \(y = \dfrac{mx - 2m - 3}{x - m}\) với \(m\) là tham số. Gọi \(S\) là tập hợp tất cả các giá trị nguyên của \(m\) để hàm số đồng biến trên các khoảng xác định. Tìm số phần tử của \(S\).}{%
    \begin{tasks}(4)
        \task Vô số
        \task 3
        \task 5
        \task 4
    \end{tasks}
}

\textit{\textbf{Dạng 2: Tìm m để hàm số nhất biến đơn điệu trên khoảng cho trước}}

\drawdashlines{33}
\newpage


% --- CÂU 4 ---
\questionLinesManual{30}{4}{Cho hàm số \(f(x) = \dfrac{mx - 4}{x - m}\) (\(m\) là tham số thực). Có bao nhiêu giá trị nguyên của \(m\) để hàm số đã cho đồng biến trên khoảng \((0; +\infty)\)?}{%
    \begin{tasks}(4)
        \task 5
        \task 4
        \task 3
        \task 2
    \end{tasks}
}

% --- CÂU 5 ---
\questionLinesManual{30}{5}{Tập hợp tất cả các giá trị thực của tham số \(m\) để hàm số \(y = \dfrac{x + 4}{x + m}\) đồng biến trên khoảng \((-\infty; -7)\) là}{%
    \begin{tasks}(4)
        \task \([4; 7)\)
        \task \((4; 7]\)
        \task \((4; 7)\)
        \task \((4; +\infty)\)
    \end{tasks}
}

% --- CÂU 6 ---
\questionLinesManual{30}{6}{Tập hợp tất cả các giá trị thực của tham số \(m\) để hàm số \(y = \dfrac{x + 5}{x + m}\) đồng biến trên khoảng \((-\infty; -8)\) là}{%
    \begin{tasks}(4)
        \task \((5; +\infty)\)
        \task \((5; 8]\)
        \task \([5; 8)\)
        \task \((5; 8)\)
    \end{tasks}
}

\textit{\textbf{Dạng 3: Tìm m để hàm số bậc 3 đơn điệu trên khoảng cho trước}}

\drawdashlines{33}
\newpage

% --- CÂU 7 ---
\questionLinesManual{30}{7}{Tập hợp tất cả các giá trị thực của tham số \(m\) để hàm số \(y = x^3 - 3x^2 + (4 - m)x\) đồng biến trên khoảng \((2; +\infty)\) là}{%
    \begin{tasks}(4)
        \task \((-\infty; 1]\)
        \task \((-\infty; 4]\)
        \task \((-\infty; 1)\)
        \task \((-\infty; 4)\)
    \end{tasks}
}

% --- CÂU 8 ---
\questionLinesManual{30}{8}{Tập hợp tất cả các giá trị của tham số \(m\) để hàm số \(y = x^3 - 3x^2 + (5 - m)x\) đồng biến trên khoảng \((2; +\infty)\) là}{%
    \begin{tasks}(4)
        \task \((-\infty; 2)\)
        \task \((-\infty; 5)\)
        \task \((-\infty; 5]\)
        \task \((-\infty; 2]\)
    \end{tasks}
}

% --- CÂU 9 ---
\questionLinesManual{30}{9}{Tập hợp tất cả các giá trị thực của tham số \(m\) để hàm số \(y = x^3 - 3x^2 + (2 - m)x\) đồng biến trên khoảng \((2; +\infty)\) là}{%
    \begin{tasks}(4)
        \task \((-\infty; -1]\)
        \task \((-\infty; 2)\)
        \task \((-\infty; -1)\)
        \task \((-\infty; 2]\)
    \end{tasks}
}

\textit{\textbf{Dạng 4: Tìm m để hàm số khác đơn điệu trên khoảng cho trước}}

\drawdashlines{33}
\newpage

% --- CÂU 10 ---
\questionLinesManual{30}{10}{Tìm tất cả các giá trị thực của tham số \(m\) sao cho hàm số \(y = \dfrac{\tan x - 2}{\tan x - m}\) đồng biến trên khoảng \(\left(0; \dfrac{\pi}{4}\right)\).}{%
    \begin{tasks}(2)
        \task \(m \leq 0\) hoặc \(1 \leq m < 2\)
        \task \(m \leq 0\)
        \task \(1 \leq m < 2\)
        \task \(m \geq 2\)
    \end{tasks}
}

% --- CÂU 11 ---
\questionLinesManual{30}{11}{Có bao nhiêu giá trị nguyên âm của tham số \(m\) để hàm số \(y = x^3 + mx - \dfrac{1}{5x^5}\) đồng biến trên khoảng \((0; +\infty)\)}{%
    \begin{tasks}(4)
        \task 0
        \task 4
        \task 5
        \task 3
    \end{tasks}
}

% --- CÂU 12 ---
\questionLinesManual{30}{12}{Gọi \(S\) là tập hợp tất cả các giá trị của tham số \(m\) để hàm số \(f(x) = \dfrac{1}{5}m^2x^5 - \dfrac{1}{3}mx^3 + 10x^2 - (m^2 - m - 20)x\) đồng biến trên \(\mathbb{R}\). Tổng giá trị của tất cả các phần tử thuộc \(S\) bằng}{%
    \begin{tasks}(4)
        \task \(\dfrac{5}{2}\)
        \task \(-2\)
        \task \(\dfrac{1}{2}\)
        \task \(\dfrac{3}{2}\)
    \end{tasks}
}


\end{document}